%! Exponential Functions — Unit 4, Lesson 4.3 — Homework
\documentclass[10pt]{article}
\usepackage{precalculus-article}
\usepackage{precalculus-boxes}
\usepackage{grading}
\gradingassignment{precalc-u04-l03-hw}
\gradingcourse{precalculus}
\newcommand{\TallMath}[1]{$\displaystyle #1\rule[-1.4em]{0pt}{3.2em}$}

\begin{document}

\pageheader{Unit 4, Lesson 4.3}{Homework}

\vspace{0.10in}

\begin{enumerate}[itemsep=16pt]

\item For each exponential function below, identify $a$, $b$, classify as growth or decay,
      and state the range.

      \begin{enumerate}[label=(\alph*), itemsep=10pt]
        \item $f(x) = 6 \cdot 2^x$

              \vspace{0.04in}
              $a = $ \smash{\gradedblank[id=hw1a1, type=short, points=1, answer={6}]{2cm}{13pt}} \quad $b = $ \smash{\gradedblank[id=hw1a2, type=short, points=1, answer={2}]{2cm}{13pt}} \quad
              Growth or decay: \smash{\gradedblank[id=hw1a3, type=short, check=text, points=1, answer={growth}]{3cm}{13pt}} \quad Range: \smash{\gradedblank[id=hw1a4, type=short, check=text, points=1, answer={(0, \infty)}, display={$(0,\infty)$}]{3cm}{13pt}}

        \item $g(x) = 12 \cdot (0.25)^x$

              \vspace{0.04in}
              $a = $ \smash{\gradedblank[id=hw1b1, type=short, points=1, answer={12}]{2cm}{13pt}} \quad $b = $ \smash{\gradedblank[id=hw1b2, type=short, points=1, answer={0.25}]{2cm}{13pt}} \quad
              Growth or decay: \smash{\gradedblank[id=hw1b3, type=short, check=text, points=1, answer={decay}]{3cm}{13pt}} \quad Range: \smash{\gradedblank[id=hw1b4, type=short, check=text, points=1, answer={(0, \infty)}, display={$(0,\infty)$}]{3cm}{13pt}}

        \item $h(x) = 0.5 \cdot 4^x$

              \vspace{0.04in}
              $a = $ \smash{\gradedblank[id=hw1c1, type=short, points=1, answer={0.5}]{2cm}{13pt}} \quad $b = $ \smash{\gradedblank[id=hw1c2, type=short, points=1, answer={4}]{2cm}{13pt}} \quad
              Growth or decay: \smash{\gradedblank[id=hw1c3, type=short, check=text, points=1, answer={growth}]{3cm}{13pt}} \quad Range: \smash{\gradedblank[id=hw1c4, type=short, check=text, points=1, answer={(0, \infty)}, display={$(0,\infty)$}]{3cm}{13pt}}
      \end{enumerate}

\item Determine whether each table represents an \textbf{exponential}, \textbf{linear}, or
      \textbf{neither} function. Show the ratios or differences that support your answer.

      \begin{enumerate}[label=(\alph*), itemsep=12pt]
        \item \renewcommand{\arraystretch}{1.3}
              \begin{tabular}{|c|c|c|c|c|}
              \hline
              $x$ & 0 & 1 & 2 & 3 \\ \hline
              $y$ & 2 & 6 & 18 & 54 \\ \hline
              \end{tabular}

              \vspace{0.04in}
              Work: \blank{6cm} \quad Type: \smash{\gradedblank[id=hw2a, type=short, check=text, points=1, answer={exponential}]{3cm}{13pt}}

        \item \renewcommand{\arraystretch}{1.3}
              \begin{tabular}{|c|c|c|c|c|}
              \hline
              $x$ & 0 & 1 & 2 & 3 \\ \hline
              $y$ & 5 & 9 & 13 & 17 \\ \hline
              \end{tabular}

              \vspace{0.04in}
              Work: \blank{6cm} \quad Type: \smash{\gradedblank[id=hw2b, type=short, check=text, points=1, answer={linear}]{3cm}{13pt}}

        \item \renewcommand{\arraystretch}{1.3}
              \begin{tabular}{|c|c|c|c|c|}
              \hline
              $x$ & 0 & 1 & 2 & 3 \\ \hline
              $y$ & 1 & 3 & 8 & 16 \\ \hline
              \end{tabular}

              \vspace{0.04in}
              Work: \blank{6cm} \quad Type: \smash{\gradedblank[id=hw2c, type=short, check=text, points=1, answer={neither}]{3cm}{13pt}}
      \end{enumerate}

\item A function $q$ has outputs 4, 20, 100, 500 for inputs 0, 1, 2, 3. Write $q(x)$
      in the form $ab^x$.

      \vspace{0.04in}
      $a = $ \smash{\gradedblank[id=hw3a, type=short, points=1, answer={4}]{2cm}{13pt}} \quad $b = $ \smash{\gradedblank[id=hw3b, type=short, points=1, answer={5}]{2cm}{13pt}} \quad $q(x) = $ \smash{\gradedblank[id=hw3c, type=short, points=1, answer={4 \cdot 5^x}]{5cm}{13pt}}

\item State the end behavior of each function using limit notation.

      \begin{enumerate}[label=(\alph*), itemsep=10pt]
        \item $f(x) = 7 \cdot 3^x$

              \vspace{0.04in}
              $\displaystyle\lim_{x\to\infty} f(x) = $ \smash{\gradedblank[id=hw4a1, type=short, points=1, answer={\infty}]{2.5cm}{13pt}} \qquad
              $\displaystyle\lim_{x\to-\infty} f(x) = $ \smash{\gradedblank[id=hw4a2, type=short, points=1, answer={0}]{2.5cm}{13pt}}

        \item $g(x) = 8 \cdot (0.6)^x$

              \vspace{0.04in}
              $\displaystyle\lim_{x\to\infty} g(x) = $ \smash{\gradedblank[id=hw4b1, type=short, points=1, answer={0}]{2.5cm}{13pt}} \qquad
              $\displaystyle\lim_{x\to-\infty} g(x) = $ \smash{\gradedblank[id=hw4b2, type=short, points=1, answer={\infty}]{2.5cm}{13pt}}
      \end{enumerate}

\item \textbf{(AP-Style Multiple Choice)} Which of the following functions has a horizontal
      asymptote at $y = 0$ \textit{and} is always decreasing on its entire domain?

      \medskip
      \begin{enumerate}[label=(\Alph*), itemsep=4pt]
        \item $f(x) = 3 \cdot 2^x$
        \item $g(x) = -4 \cdot (0.5)^x$
        \item $h(x) = 5 \cdot (0.8)^x$
        \item $k(x) = 2 \cdot 1^x$
      \end{enumerate}

      \vspace{0.04in}
      \leavevmode\rlap{\textit{Letter:}}\phantom{Answer:}\spacefactor2000{} \smash{\gradedblank[id=hw5a, type=choice, points=1, answer=C]{1.5cm}{13pt}} \quad Justify your choice in one sentence:

      \vspace{0.04in}
      \writebox{\raisebox{-0.75pt}[0pt][0pt]{\gradedblank[id=hw5b, type=rubric, points=2,
  display={$h(x) = 5\cdot(0.8)^x$: $a = 5 > 0$ and $0 < b = 0.8 < 1$, so it always decreases toward its asymptote $y = 0$.}]{\linewidth}{15.5pt}}}

\item \textbf{(Extension)} The definition of an exponential function requires $b \ne 1$.
      Explain in two or more sentences why the function $f(x) = a \cdot 1^x$ is excluded.
      What kind of function is it, and what property of exponential functions does it fail to satisfy?

      \vspace{0.04in}
      \writegap\\
      \writegap\\
      \writebox{\gradedblank[id=hw6, type=rubric, points=2,
  display={$f(x) = a\cdot 1^x = a$ for all $x$: a constant function. Its outputs never change, so it is neither increasing nor decreasing (no proportional growth or decay) and its limit in both directions is simply $a$.}]{\linewidth}{40pt}}\\

\end{enumerate}

\end{document}
